\documentclass{article}
\title{A local teaching example of mean and range}
\author{PaperSpine documentation example}
\date{}
\begin{document}
\maketitle
\begin{abstract}
This teaching note uses five deliberately constructed values to demonstrate
the arithmetic mean and range. It is not an empirical research result.
\end{abstract}
\section{Purpose and materials}
The accompanying measurements.csv contains A=2, B=4, C=6, D=8 and E=10.
The purpose is to explain how the mean summarizes these five numbers while
the range retains information about their spread. No participants, observations,
external sources or claims of novelty are involved.
\section{Calculation}
The sum is 2+4+6+8+10=30. Dividing by the count of five gives a mean of 6.
The minimum is 2 and maximum is 10, so the range is 10-2=8.
\section{Figure and interpretation}
The editable figure values.svg shows five bars with heights 2, 4, 6, 8 and 10.
Each bar is labelled with its value. The horizontal axis labels are identifiers,
not time points. The mean alone does not show the minimum or maximum.
\section{Limitations}
These values are constructed solely for documentation. No uncertainty estimate,
population inference, journal suitability or submission readiness is claimed.
\section{Local verification}
A reader can verify the CSV rows, the sum and count, and the five bar labels.
This note and its editable figure may be copied, modified and redistributed.
\end{document}
